
The majority of quantization schemes focus on compressing LLMs by using \textit{weight-only quantization}, \citep{gptq, spqr, awq, AQLM, QuIPsharp}. These methods downcast each weight into a low-precision representation and upcast it before the actual computation. The main computation is still performed in high precision. 
Several works show that, unlike weights, quantizing the activations is hard due to the outlier features \citep{outlier_suppression, llm_int8, smoothquant}. For 8-bit case, LLM.int8() \citep{llm_int8} identifies the outlier features during inference and keeps them in 16 bits which results in poor performance. SmoothQuant \citep{smoothquant} normalizes the features using some scaling factors from a calibration set, solving the issue for the 8-bit case at the cost of introducing extra hyper-parameters. For 4-bit quantization, recent studies identify the outlier features offline and keep them in high precision. Atom \citep{atom} developed a complex kernel for mixed-precision MatMul in the presence of outliers while QUIK \citep{quik} keeps the down-projection layer in 8 bits.


Two weight-only quantization methods, QuIP \citep{QuIP} and QuIP\# \citep{QuIPsharp} have previously considered improving quantization by applying rotations.  \citet{QuIP} introduced the idea of \emph{incoherence processing} which applies rotation matrices to the left and right of each weight matrix, as well as the Hessian, which is used in minimizing the weight-quantization objective. \citet{xi2023training} uses a similar idea during training, using exact Hadamard transformations for each linear layer in the forward pass.


Finally, KV cache quantization is another line of research that aims to compress the cached keys and values during the generation phase. This is crucial for large batch size and long-context length generation as the KV cache will be the main memory bottleneck in such problems. \citet{sheng2023flexgen} quantizes the KV cache using 4-bit group-wise quantization. KVQuant \citep{kvquant} pushes this limit to 3-bit quantization and KIVI \citep{kivi} shows promising results on 2-bit KV cache quantization. Such methods show that outliers also exist in the keys, and apply a set of complex ideas (like feature-wise quantization, non-uniform representation, and keeping high precision outliers) to recover the accuracy of a quantized KV cache.


In this work we also adopt the Hadamard transform to improve quantization of weights through incoherence processing. Instead of undoing the Hadamard transform during the forward pass, we adopt the computational invariance theorem from SliceGPT \citep{slicegpt} to fuse the transformations into the weights where possible. Instead of requiring two Hadamard transforms per weight-matrix in the forward pass, QuaRot requires just $1\tfrac{1}{2}$ Hadamard transforms per transformer layer. Computational invariance also means that the \emph{activations} are incoherence-processed, enabling them to be effectively quantized. We also apply a similar technique to the attention block and quantize the KV cache in 4 bits with minimal accuracy loss.
